\documentclass[UTF8,a4paper,11pt]{ctexart}
\usepackage[a4paper,margin=2.3cm]{geometry}
\usepackage{amsmath,amssymb,mathtools}
\usepackage{enumitem}
\usepackage[hidelinks]{hyperref}
\usepackage{fancyhdr}
\setlist[itemize]{leftmargin=2.2em,itemsep=0.25em,topsep=0.35em}
\setlength{\parindent}{2em}\setlength{\parskip}{0.35em}
\pagestyle{fancy}\fancyhf{}\fancyhead[L]{离散数学与结构（I）}\fancyhead[R]{Lec 09}\fancyfoot[C]{\thepage}\setlength{\headheight}{14pt}
\title{环、理想与商环}\author{zhangjunshuo}\date{2026 年 10 月}
\begin{document}
\maketitle
\begin{quote}
本讲根据 2026 年 10 月 13 日课堂转写、手写原稿和课程教材整理。不同教材对环是否含乘法单位元有不同约定，正文明确采用含幺环作为默认语境；需要不含幺版本时会特别说明。

课堂的讲法是“走马观花过概念”，不少标准结论只给一句话或留作习题。本整理版按教材补全了这些地方，并在每节末尾标出出处：有限整环是域（第 1 节）、第二/第三同构定理与 $M_n(F)$ 无平凡理想（第 2 节）、中国剩余定理的满射证明与正交幂等元（第 3 节）、$R[x]$ 单位的完整分层与素子环（第 4 节）。
\end{quote}
\tableofcontents\clearpage

\section*{记号速查}
\begin{itemize}
\item $R,S$：环；$1_R$：单位元；$R^\times$：单位群；$R[x]$：多项式环。
\item $I,J$：理想；$(a)$：主理想；$R/I$：商环；$\operatorname{Frac}(R)$：交换整环的分式域。
\item $\operatorname{char}(R)$：特征；$\mathbb F_p=\mathbb Z/p\mathbb Z$：$p$ 元素的有限域。
\end{itemize}

\section*{勘误提醒}
课堂转写与教材核对后，有六处结论需要补上前提或改正表述，正文都已写清：
\begin{enumerate}
\item \textbf{“$F[x]$ 中常数项非零就可逆”只在 $F$ 是域时成立。} 完整分层是：域上 $F[x]^\times=F\setminus\{0\}$；整环上 $R[x]^\times=R^\times$；一般交换含幺环上要求 $a_0$ 可逆且 $a_1,\dots,a_n$ 全为幂零元。反例：在 $(\mathbb Z/4\mathbb Z)[x]$ 中 $(1+2x)^2=1$，所以 $1+2x$ 可逆。
\item \textbf{“特征必为素数”只对整环成立。} $\mathbb Z/4\mathbb Z$ 的特征是 $4$。
\item \textbf{$I\cap J=IJ$ 需要 $I+J=R$，并且用到单位元。} 在 $R=2\mathbb Z$、$I=J=2\mathbb Z$ 时 $I\cap J=2\mathbb Z$ 而 $IJ=4\mathbb Z$。
\item \textbf{“每个极大理想都是素理想”需要 $R$ 交换。} 非交换反例：$M_n(F)$（$n\ge2$）没有非平凡理想，故 $(0)$ 极大；但 $E_{11}E_{22}=0$，所以 $(0)$ 不素。
\item \textbf{整环的定义不能省“含幺且 $1\ne0$”。} 否则零环 $\{0\}$ 也会成为整环。
\item \textbf{$(a)(b)=(ab)$ 不是 $\mathbb Z$ 独有}，任何交换环的主理想都满足；不要与一般的 $IJ\subseteq I\cap J$ 混为一谈。
\end{enumerate}
另外两处约定问题：本讲默认环含单位元、同态保持单位元，因此子环判定要求 $1\in S$（故含幺环中理想 $I$ 是子环当且仅当 $I=R$）；改用不含幺约定（[DF] \S7.1）时 $2\mathbb Z$ 又成了 $\mathbb Z$ 的子环，两种约定都自洽，混用才会错。
\clearpage
\section{环的定义与典型例子}


\textbf{本节用到的记号，先一次交代清楚}（更早的讲次已逐个定义过，这里只复述结论）：

\begin{itemize}
\item $\mathbb Z,\mathbb Q,\mathbb R,\mathbb C$：整数、有理数、实数、复数；$\mathbb H$ 是四元数环，定义见 1.4 小节。
\item $\mathbb Z/n\mathbb Z$：整数模 $n$ 的剩余类环，元素是 $0,1,\dots,n-1$ 所在的等价类，加减乘都在模 $n$ 下进行。
\item $\mathbb F_p=\mathbb Z/p\mathbb Z$（$p$ 为素数）：含 $p$ 个元素的有限域。特别地 $\mathbb F_2=\mathbb Z/2\mathbb Z=\{0,1\}$，其中 $1+1=0$。
\item $R[x]$：系数在 $R$ 中的一元多项式环；$\sum_i$ 是有限求和，$\prod_i$ 是有限连乘。
\item $E_{ij}$：\textbf{矩阵单位}，即第 $i$ 行第 $j$ 列是 $1$、其余位置是 $0$ 的矩阵（1.3 小节举例时用到）。
\item $S\subseteq R$：$S$ 是 $R$ 的子集；$R^S$：从集合 $S$ 到环 $R$ 的全体函数。
\item $\ker f$ 与 $\operatorname{im}f$：同态 $f$ 的核与像，定义见 1.5 小节。
\end{itemize}


\subsection{环的两种约定}


环 $R$ 是一个集合，带有加法与乘法，使 $(R,+)$ 为阿贝尔群，乘法满足结合律，并满足左右分配律

\[a(b+c)=ab+ac,\qquad (a+b)c=ac+bc.\]

有的教材把乘法单位元 $1$ 写入定义，有的教材允许没有单位元。本文默认讨论\textbf{含幺环}；若只需环的加法群、乘法结合律和分配律，则称为不必含幺的环。若乘法还交换，称为交换环。

非零元 $a,b$ 满足 $ab=0$ 时，称它们是零因子；\textbf{交换含幺环且 $1\ne0$、没有非零零因子}时称为整环（integral domain）。这里“含幺且 $1\ne0$”不能省：否则零环 $\{0\}$ 也会被算作整环，而 $\{0\}$ 显然不是我们想要的对象。除环（division ring，也称斜域）要求每个非零元都有乘法逆元；交换的除环就是域。注意“除环”只去掉乘法交换，不放弃单位元；四元数环 $\mathbb H=\{a+bi+cj+dk:a,b,c,d\in\mathbb R\}$ 就是非交换除环的标准例子（$i,j,k$ 满足 $i^2=j^2=k^2=ijk=-1$，详见 1.4 小节）。

整环与域之间有一个只用有限性就能跨过的台阶：

\textbf{命题。} 有限整环必为域。

\textbf{证明。} 设 $R$ 是有限整环，取 $a\ne0$。考虑映射 $L_a:R\to R$，$x\mapsto ax$。若 $ax=ay$，由消去律得 $x=y$，所以 $L_a$ 是单射。$R$ 有限，单射即满射，故存在 $x\in R$ 使 $ax=1$。这说明每个非零元 $a$ 都有右逆；同样的论证（用 $x\mapsto xa$）给出左逆，两者相等，所以 $a$ 是单位。于是 $R$ 是域。$\square$

这条命题解释了为什么 $\mathbb Z/n\mathbb Z$ 在 $n$ 为素数时是域：它本身是有限整环。也说明“域 = 整环 + 每个非零元可逆”这个条件在有限情形下自动满足。

\subsection{零因子、消去律与环的层级}


课堂把零因子写在整环的定义旁边，这两件事其实是一体的：

\textbf{命题。} 交换含幺环 $R$ 是整环，当且仅当 $R$ 中消去律成立：对 $a\ne0$，由 $ab=ac$ 可得 $b=c$。

\textbf{证明。} 若 $R$ 有非零零因子，取 $a\ne0$ 与 $b\ne0$ 使 $ab=0=a\cdot0$，则 $ab=a\cdot0$ 但 $b\ne0$，消去律失效。反之若消去律失效，即存在 $a\ne0$ 与 $b\ne c$ 使 $ab=ac$，则 $a(b-c)=0$ 且 $b-c\ne0$，于是 $a$ 是非零零因子。$\square$

非交换时还要区分左右：$ab=0$（$a,b\ne0$）时 $a$ 叫左零因子、$b$ 叫右零因子；$\mathbb Z/4\mathbb Z$ 中的 $2$ 同时是左右零因子（$2\cdot2=0$）。

把这几层关系记成一条链会清楚很多：


\begin{center}
\begin{tabular}{ll}
环（不一定含幺） & 加法阿贝尔群 $+$ 乘法半群 $+$ 分配律\\ \hline
含幺环 & 再加 $1\ne0$，$1a=a1=a$\\
交换环 & 乘法交换\\
整环 & 交换含幺 $+$ 无零因子\\
除环 & 每个非零元有逆（乘法可不交换）\\
域 & 交换的除环 \;=\; 整环 $+$ 非零元全可逆
\end{tabular}
\end{center}


包含关系是：域 $\subset$ 除环、域 $\subset$ 整环，而 $\mathbb H$ 是除环但不是域，$\mathbb Z$ 是整环但不是域。两条容易搞混的边界：\textbf{除环不要求乘法交换}（$\mathbb H$），\textbf{整环不要求每个非零元可逆}（$\mathbb Z$）。

\subsection{例子}


\begin{itemize}
\item $\mathbb Z,\mathbb Q,\mathbb R,\mathbb C$ 都是含幺交换环；$\mathbb Z$ 是整环但不是域。
\item $\mathbb Z/n\mathbb Z$ 是交换环；当且仅当 $n$ 为素数时它是域。
\item 若 $R$ 是环，则 $M_n(R)$ 在逐项加法和矩阵乘法下是环。只要 $R$ 不是零环且 $n\ge2$，$M_n(R)$ \textbf{必}不交换——例如 $E_{12}E_{21}=E_{11}$ 而 $E_{21}E_{12}=E_{22}$，两者不等。当 $F$ 是域且 $n\ge2$ 时，$M_n(F)$ 还没有非平凡的\textbf{双边理想}（见第 2 节），所以它不是域也不是除环。
\item 阿贝尔群 $A$ 的自同态集合 $\operatorname{End}(A)$ 在逐点相加、复合下是环；复合一般不交换。
\item 环的有限直积 $R_1\times\cdots\times R_k$ 按坐标运算成为环。
\item 从集合 $S$ 到环 $R$ 的所有函数 $R^S$ 按逐点运算成为环；有限 $S$ 时它就是有限直积。
\end{itemize}


\subsection{多项式环与群环}


若 $R$ 为交换含幺环，形式多项式
\[R[x]=\left\{\sum_{i=0}^n a_i x^i:a_i\in R\right\}\]
按系数相加、卷积相乘构成交换环。多项式中的 $x$ 是形式符号，不必先赋予数值。

更一般地，设 $R$ 是交换含幺环、$G$ 是群（不要求交换），群环 $R[G]$ 的元素是有限支撑形式和 $\sum_{g\in G}a_g g$，乘法由群乘法双线性展开：

\[
\Bigl(\sum_{g}a_g g\Bigr)\Bigl(\sum_{h}b_h h\Bigr)=\sum_{g,h}a_gb_h\,(gh).
\]

判断交换性有一个干净的判据：\textbf{$R[G]$ 交换当且仅当 $G$ 交换。} 因为基元 $g,h$ 在 $R[G]$ 中的乘积就是群乘积 $gh$，所以两组基元两两交换等价于 $G$ 交换；反过来 $G$ 交换时把两个形式和的乘积展开即可见其对称。特别地，$\mathbb F_2[S_3]$ 非交换。当 $G$ 为无限群时仍只允许有限项（有限支撑），否则乘积无定义。

四元数环 $\mathbb H=\{a+bi+cj+dk:a,b,c,d\in\mathbb R\}$ 满足
\[i^2=j^2=k^2=ijk=-1,\]
例如 $ij=k$ 而 $ji=-k$，所以它是非交换除环。非零四元数 $q=a+bi+cj+dk$ 的逆为
\[q^{-1}=\frac{a-bi-cj-dk}{a^2+b^2+c^2+d^2}.\]

\subsection{子环、同态与核}


子集 $S\subseteq R$ 若在同一加法、乘法下构成环，则称为子环。\textbf{这里有一个约定分歧，必须说清}：

\begin{itemize}
\item 若“环”一律指含幺环（本讲采用的约定，也是 [DS] 的约定），子环判定是：\textbf{$1_R\in S$，且 $S$ 对减法与乘法封闭}。此时 $2\mathbb Z\subseteq\mathbb Z$ 不是子环，因为它不含 $1$。
\item 若允许不含幺的环（[DF] §7.1 的约定），子环只需是“对乘法封闭的加法子群”，于是 $2\mathbb Z$ 是 $\mathbb Z$ 的子环。
\end{itemize}


两种约定都自洽，混用才会出错。采用含幺约定时有一个直接推论：\textbf{含幺环 $R$ 的理想 $I$ 是子环当且仅当 $I=R$}，因为 $1\in I$ 会迫使 $I=R$。这也解释了为什么“$\ker f$ 是子环”这句话在含幺语境下要改说成“$\ker f$ 是理想”（见第 2 节）。

环同态 $f:R\to S$ 满足 $f(a+b)=f(a)+f(b)$、$f(ab)=f(a)f(b)$；在含幺约定下还要求 $f(1_R)=1_S$，这样 $\operatorname{im}f$ 才是 $S$ 的（含幺）子环。

同态的核
\[\ker f=\{a\in R:f(a)=0\}\]
是 $R$ 的理想；像 $\operatorname{im}f$ 是 $S$ 的子环。这为下一节的商环与同构定理作准备。

\textbf{关于单位元的约定要小心。} 若采用“环不必含幺”的约定，则 $f(1_R)$ 未必是 $S$ 的单位元：例如 $\mathbb Z\to\mathbb Z\times\mathbb Z$、$n\mapsto(n,0)$ 保持加法与乘法，但把 $1$ 送到 $(1,0)$，而 $(1,0)$ 在 $\mathbb Z\times\mathbb Z$ 中不是单位元。所以在含幺环的语境里，通常\textbf{额外要求}同态满足 $f(1_R)=1_S$；本讲采用这个约定，凡不要求保幺的地方会单独说明。

\subsection{参考对照}


\begin{itemize}
\item [DF] §7.1（环的定义与例子，含 $M_n(R)$、$\mathbb H$、群环）、§7.2（多项式环、整环、$R[x]^\times=R^\times$）。
\item [DS] §10.1（环和域的定义、整环、四元数体、有限整环是域）。
\item 原始材料：$../../../课堂手记和录音转写/Lec09/$。
\end{itemize}

\section{理想与商环}


\textbf{本节用到的记号，先一次交代清楚：}

\begin{itemize}
\item $I\trianglelefteq R$：$I$ 是环 $R$ 的理想（符号借用群论的正规子群，含义由下文定义给出）。$I\le R$ 表示 $I$ 是子环。
\item $n\mathbb Z=\{nk:k\in\mathbb Z\}$：$n$ 的全体倍数，是 $\mathbb Z$ 的理想。
\item $(X)$：由子集 $X$ 生成的理想；$(a)$ 是单元素生成的主理想。
\item $\gcd(a,b)$、$\operatorname{lcm}(a,b)$：最大公因数与最小公倍数，只对整数（或更一般的唯一分解整环）使用。
\item $\operatorname{diag}(d_1,\dots,d_n)$：以 $d_1,\dots,d_n$ 为对角元、其余位置为 $0$ 的 $n\times n$ 对角矩阵。
\item $R[x]$：系数在 $R$ 中的多项式环；$\deg f$ 是多项式 $f$ 的次数；\textbf{首一}指最高次项系数为 $1$。
\end{itemize}


\subsection{理想}


环 $R$ 的加法子群 $I$ 称为（双边）理想，若对任意 $r\in R$、$a\in I$ 都有
\[ra\in I,\qquad ar\in I.\]
在交换环中左右条件相同。理想是环同态核的正确替代物；任意子环未必能作为商环的陪集。

\textbf{理想与子环的关系要分清。} 由定义，理想对减法封闭、对自身乘法封闭，所以它是“非含幺意义下”的子环；但它还额外要求吸收整个环的乘法。在含幺环里这个差别立刻显现：\textbf{理想 $I\ne R$ 不可能包含 $1_R$}，因为一旦 $1\in I$ 就有 $r=r\cdot1\in I$ 对所有 $r$ 成立，于是 $I=R$。

所以同态的核是理想，而且一般\textbf{不是}含幺子环——例如 $\mathbb Z\to\mathbb Z/2\mathbb Z$ 的核 $2\mathbb Z$ 不含 $1$。课堂手写稿里写的“$\ker f$ 是子环”，在含幺约定下应改述为“$\ker f$ 是理想，因而只是非含幺意义下的子环”。反过来，子环 $A\le R$ 成为理想，当且仅当它吸收乘法：$rA\subseteq A$ 且 $Ar\subseteq A$ 对所有 $r\in R$。

\textbf{理想的基本运算。} 设 $I,J\trianglelefteq R$。和与交定义为

\[
I+J=\{a+b:a\in I,\ b\in J\},\qquad I\cap J=\{x:x\in I\text{ 且 }x\in J\},
\]

两者都是理想：$I+J$ 还是包含 $I\cup J$ 的最小理想（即 $I\cup J$ 生成的理想），$I\cap J$ 是同时包含于两者的最大理想。

\textbf{理想的乘积}是第三个运算，定义为

\[
IJ=\left\{\sum_{k=1}^n a_kb_k:a_k\in I,\ b_k\in J,\ n\ge1\right\},
\]

即“$I$ 的元素乘 $J$ 的元素，再取有限和”所成的集合。它是包含所有乘积 $ab$（$a\in I,b\in J$）的最小理想；不能把它误写成只含单个乘积的集合。

三个运算之间有包含链

\[
IJ\;\subseteq\;I\cap J\;\subseteq\;I+J\;\subseteq\;R,
\]

其中 $IJ\subseteq I\cap J$ 是因为每个乘积 $ab$ 既在 $I$ 又在 $J$；$I\cap J\subseteq I+J$ 是因为 $x=x+0$。此外成立分配律

\[
I(J+K)=IJ+IK,\qquad (I+J)K=IK+JK,
\]

并可以定义幂 $I^n=\underbrace{I\cdots I}_{n\text{ 个}}$。

注意一般的 $I\cup J$ \textbf{不是}理想：例如在 $\mathbb Z$ 中 $2\in(2)$、$3\in(3)$，但 $2+3=5\notin(2)\cup(3)$，对加法不封闭。这正是需要引入 $I+J$ 的原因。

理想的交仍是理想（任意多个理想的交仍是理想）。由子集 $X\subseteq R$ 生成的理想记为 $(X)$，是包含 $X$ 的最小理想。对单元素 $a$，$(a)$ 称主理想；在交换含幺环中
\[ (a)=\{ra:r\in R\}.\]
一般的双边生成理想由有限和
\[\sum_i r_i x_i s_i\]
组成；交换环中可写成 $\sum_i r_i x_i$。这里 $r_i,s_i$ 遍历整个 $R$，$x_i$ 遍历 $X$，项数有限。之所以必须写成这种“两侧都乘”的形式：非交换时 $ra$ 与 $as$ 未必相同，而理想要求对左右乘法都封闭。验证也不难——两个这样的有限和相加仍是有限和，再左乘或右乘一个 $r\in R$ 只是把每项变成 $r(r_ix_is_i)= (rr_i)x_is_i$，形式不变，所以这个集合确实是对左右乘法封闭的加法子群，因而是理想；又因为每一项都必须落在任何包含 $X$ 的理想里，它是最小的。

在 $\mathbb Z$ 中每个理想都是主理想：$I=(d)=d\mathbb Z$，其中 $d$ 是 $I$ 中非零元素的最大公因数。于是

\[ (a)\cap(b)=(\operatorname{lcm}(a,b)),\qquad (a)+(b)=(\gcd(a,b)), \]

这里 $\gcd(a,b)$ 是 $a,b$ 的最大公因数、$\operatorname{lcm}(a,b)$ 是最小公倍数。第一条是因为“同时是 $a$ 和 $b$ 的倍数”就是“$\operatorname{lcm}(a,b)$ 的倍数”，第二条是因为 $\mathbb Z$ 中能同时由 $a,b$ 线性表出的数恰是 $\gcd(a,b)$ 的倍数。

回到乘积 $IJ$：对\textbf{主理想}，它直接由生成元给出。在交换环中

\[
(a)(b)=(ab),\qquad (a_1)(a_2)\cdots(a_k)=(a_1a_2\cdots a_k).
\]

这不是 $\mathbb Z$ 才有的性质，任何交换环都成立（[DF] §7.6）。容易和它混淆的是另一个式子：一般理想只有 $IJ\subseteq I\cap J$（每一项 $ab$ 既在 $I$ 又在 $J$），反向包含需要 $I+J=R$，见第 3 节。

\subsection{商环的构造}


理想 $I\trianglelefteq R$ 给出加法陪集
\[R/I=\{a+I:a\in R\}.\]
定义
\[ (a+I)+(b+I)=(a+b)+I,\qquad (a+I)(b+I)=ab+I.\]
乘法良定义的关键是：若 $a'=a+i$、$b'=b+j$，则
\[a'b'-ab=aj+ib+ij\in I.\]
因此 $R/I$ 成为环，零元为 $I$，单位元为 $1+I$（若 $R$ 含幺）。

\subsection{同构定理}


对环同态 $f:R\to S$，定义
\[\bar f:R/\ker f\longrightarrow\operatorname{im}f,\qquad a+\ker f\longmapsto f(a).\]
它良定义、双射且保持运算，因此
\[R/\ker f\cong\operatorname{im}f.\]

特别地，对理想 $I$ 的自然投影 $\pi:R\to R/I$ 的核正是 $I$。若 $f:R\to S$ 为满射，则 $R/\ker f\cong S$。

\subsection{第二、第三同构定理}


第一同构定理只处理一个同态。群论里还有两条同构定理，环论中逐字平行，证法也是把第一同构定理套到合适的映射上。课程教材把它们留作习题，这里补全陈述。

\textbf{第二同构定理。} 设 $A\le R$ 是子环，$I\trianglelefteq R$ 是理想。则 $A\cap I$ 是 $A$ 的理想、$A+I$ 是 $R$ 的子环，并且

\[
A/(A\cap I)\;\cong\;(A+I)/I.
\]

\textbf{证明提纲。} 考虑复合映射 $A\hookrightarrow A+I\twoheadrightarrow(A+I)/I$。它是满射，因为 $A+I$ 中每个元素 $a+i$（$a\in A$、$i\in I$）的像都是 $a+I$。它的核是 $\{a\in A:a\in I\}=A\cap I$。对 $A\to(A+I)/I$ 用第一同构定理即得。$\square$

\textbf{第三同构定理。} 设 $I\subseteq J$ 都是 $R$ 的理想。则 $J/I$ 是 $R/I$ 的理想，并且

\[
(R/I)\big/(J/I)\;\cong\;R/J.
\]

\textbf{证明提纲。} 考虑复合映射 $R\twoheadrightarrow R/I\twoheadrightarrow(R/I)/(J/I)$。它满，且核是 $\{r\in R:r+I\in J/I\}=J$。再用第一同构定理。$\square$

读这两条时建议回到 [Lec 06](../lec06/lec06.md) 的群版本对照：结构完全一样，只是“正规子群”换成了“理想”，而“子群”换成了“子环”。

\subsection{例：$M_n(F)$ 没有非平凡理想}


设 $F$ 是域、$n\ge2$，则矩阵环 $M_n(F)$ 只有零理想和自身。这是一个“非交换环里理想不一定多”的典型例子，也是“$M_n(F)$ 不是域”的根本原因：域要求没有非平凡理想，而 $M_n(F)$ 的非零元大多不可逆，却并不能靠理想把它切开。

\textbf{证明提纲。} 取非零理想 $I\ne(0)$，选 $0\ne A\in I$，则 $A$ 有某个非零元 $a_{ij}\ne0$。用矩阵单位 $E_{kl}$（只有一个位置为 $1$ 的矩阵）左乘、右乘 $A$，可以把 $A$ 的单个非零元“搬运”到任意位置：

\[
E_{ki}\,A\,E_{jl}=a_{ij}E_{kl}.
\]

于是 $a_{ij}E_{kl}\in I$。再用对角矩阵 $\operatorname{diag}(a_{ij}^{-1},1,\dots,1)$ 乘一次，得到 $E_{kl}\in I$。所有矩阵单位都在 $I$ 中，而它们线性张成整个 $M_n(F)$，所以 $I=M_n(F)$。$\square$

\subsection{例：求值同态与 $(x-c)$}


设 $R$ 是交换含幺环，固定 $c\in R$。求值映射

\[
\operatorname{ev}_c:R[x]\longrightarrow R,\qquad f\longmapsto f(c)
\]

是满环同态：它保持加法和乘法，且 $\operatorname{ev}_c(1)=1$。它的核是所有以 $c$ 为根的多项式。

\[
\ker\operatorname{ev}_c=(x-c).
\]

包含关系 $(x-c)\subseteq\ker\operatorname{ev}_c$ 是显然的。反向包含用\textbf{带余除法}（多项式除法：用 $g$ 除 $f$ 得到商 $q$ 与余式 $r$，且 $\deg r<\deg g$；当 $g$ 的首项系数是 $1$、即 $g$ 为\textbf{首一}多项式时，$q,r$ 一定存在且唯一）。在 $R[x]$ 中除以首一多项式 $x-c$，得

\[
f(x)=q(x)(x-c)+r,\qquad r\in R,
\]

若 $f(c)=0$ 则 $r=f(c)=0$，所以 $f\in(x-c)$。于是第一同构定理给出

\[
R[x]/(x-c)\;\cong\;R .
\]

这是“把 $R[x]$ 中 $x$ 的取值钉在 $c$ 上”的严格说法，也是后面构造域扩张（例如 $\mathbb R[x]/(x^2+1)\cong\mathbb C$）的模板。

\subsection{理想对应定理}


设 $I\trianglelefteq R$，$\pi:R\to R/I$ 是自然投影。则

\[
\{J\trianglelefteq R:I\subseteq J\}\;\longleftrightarrow\;\{K\trianglelefteq R/I\},
\qquad J\longmapsto J/I,\qquad \pi^{-1}(K)\longleftarrow K
\]

是一一对应。它保持包含关系（$J_1\subseteq J_2$ 当且仅当 $J_1/I\subseteq J_2/I$），并把和与交对应到和与交：

\[
(J_1+J_2)/I=(J_1/I)+(J_2/I),\qquad (J_1\cap J_2)/I=(J_1/I)\cap(J_2/I).
\]

注意乘积\textbf{不}被简单保持：一般只有 $(J_1/I)(J_2/I)=(J_1J_2+I)/I$，分母上多出的 $I$ 不能省。

这是群的对应定理在环论中的平行版本。下一节要用它的一个特例：$M\trianglelefteq R$ 是极大理想，当且仅当商环 $R/M$ 只有零理想和自身，也就是 $R/M$ 是域。

\subsection{参考对照}


\begin{itemize}
\item [DF] §7.3（同态、理想、商环与第一同构定理）、§7.4（理想的和、积、商与对应定理）；第二、第三同构定理的平行版本在群论 §3.3。
\item [DS] §10.2（环同态与理想）；该书把环的第二、第三同构定理留作习题。
\item 原始材料：$../../../课堂手记和录音转写/Lec09/$。
\end{itemize}

\section{极大理想、素理想与中国剩余定理}


\subsection{极大理想与域}


真理想 $M\subsetneq R$ 称为极大理想，若不存在严格夹在 $M$ 与 $R$ 之间的理想。若 $R$ 是交换含幺环，则
\[M\text{ 极大}\iff R/M\text{ 是域}.\]

证明用理想对应定理：$R/M$ 的理想恰对应包含 $M$ 的理想；商环只有零理想和自身，当且仅当 $M$ 极大。再用一条等价刻画：\textbf{交换含幺环没有非平凡理想，当且仅当它是域}（“含幺”不能省——在 $\mathbb Z/2\mathbb Z$ 的加法群上取零乘法，得到的环没有非平凡理想，却没有单位元，不是域）。域的非零元都可逆，反过来若非零元 $a$ 生成的理想只能是整个环，则 $1=ra$，所以 $a$ 可逆。

\subsection{素理想与整环}


真理想 $P$ 称为素理想，若
\[ab\in P\Longrightarrow a\in P\text{ 或 }b\in P.\]
在交换环中
\[P\text{ 素}\iff R/P\text{ 是整环}.\]
商环中 $(a+P)(b+P)=P$ 恰等价于 $ab\in P$，这就是上面等价性的计算。

\textbf{极大必素，但需要交换性。} 设 $R$ 是\textbf{交换}环、$M$ 极大，则 $R/M$ 是域，域是整环，所以 $M$ 素。非交换时结论不成立：$M_n(F)$（$n\ge2$，$F$ 是域）没有非平凡理想（见第 2 节），所以 $(0)$ 是极大理想；但 $E_{11}E_{22}=0$ 而 $E_{11},E_{22}\ne0$，说明 $M_n(F)$ 有零因子、$(0)$ 不素。

顺带说明：非交换环里“素理想”的标准定义不是 $ab\in P\Rightarrow a\in P$ 或 $b\in P$，而是对任意理想 $A,B$：

\[
AB\subseteq P\ \Longrightarrow\ A\subseteq P\ \text{或}\ B\subseteq P .
\]

两个定义在交换环中一致，所以本讲其余部分只用交换版本。

例：在 $\mathbb Z$ 中，$(p)$ 在 $p$ 为素数时既是极大理想又是素理想；$(0)$ 是素理想，因为 $\mathbb Z$ 是整环，但不是极大理想（它被任何 $(p)$ 严格包含）。

三条需要记住的相互关系：

\begin{itemize}
\item \textbf{极大必素}：若 $M$ 极大，则 $R/M$ 是域，域是整环，所以 $M$ 素。反之不然，$(0)\subsetneq\mathbb Z$ 就是反例。
\item \textbf{素不必极大}：在 $\mathbb Z$ 中只有 $(0)$ 素而不极大；在 $\mathbb Z[x]$ 中 $(x)$ 是素理想（$\mathbb Z[x]/(x)\cong\mathbb Z$ 是整环）但不极大（$\mathbb Z[x]/(x)\cong\mathbb Z$ 不是域）。可见“素”比“极大”严格更弱。
\item \textbf{商环刻画可以互推}：由理想对应定理，$R/M$ 没有非平凡理想当且仅当 $M$ 与 $R$ 之间没有别的理想；这正是极大的定义。
\end{itemize}


\subsection{互素理想与中国剩余定理}


交换环中理想 $I,J$ 互素（或互补）是指
\[I+J=R.\]
此时存在 $i\in I,j\in J$ 使 $i+j=1$，并且
\[I\cap J=IJ.\]
包含关系显然给出 $IJ\subseteq I\cap J$；反向包含由 $a=a(i+j)=ai+aj$ 得到，其中 $ai\in IJ$、$aj\in IJ$。

若 $I,J$ 互素，则中国剩余同构为

\[
R/(I\cap J)\;\longrightarrow\;R/I\times R/J,\qquad r+(I\cap J)\longmapsto(r+I,\ r+J),
\]

右端是环同态（按坐标运算），而\textbf{同态本身并不需要互素就能定义}；互素的作用是让它成为同构。这一点值得拆开看，因为“为什么需要 $I+J=R$”是本节最容易含糊的地方。

\textbf{核即为 $I\cap J$。} $r$ 映到 $(0,0)$ 当且仅当 $r\in I$ 且 $r\in J$。这一步不需要互素。

\textbf{互素恰好给出满射。} 若 $I+J=R$，取 $i\in I$、$j\in J$ 使 $i+j=1$。令

\[
e_1:=j=1-i\in J,\qquad e_2:=i=1-j\in I .
\]

逐个核对同余：$e_1=1-i$ 满足 $e_1\equiv1\pmod I$（因为 $i\in I$），又 $e_1=j\in J$ 给出 $e_1\equiv0\pmod J$；$e_2$ 对称地满足 $e_2\equiv0\pmod I$、$e_2\equiv1\pmod J$。给定任意 $(a+I,b+J)$，元素

\[
r=a\,e_1+b\,e_2
\]

就映到 $(a+I,\,b+J)$：因为 $a\,e_1\equiv a\pmod I$ 而 $b\,e_2\equiv0\pmod I$，在模 $J$ 一侧同理。所以映射满。反过来若映射满，则存在 $e_1\in R$ 映到 $(1+I,0+J)$，于是 $e_1\in J$ 且 $1-e_1\in I$，得 $1\in I+J$，即 $I+J=R$。

顺带得到一个有用的副产品：$e_1,e_2$ 是\textbf{正交幂等元}，满足

\[
e_1+e_2=1,\qquad e_1e_2=0,\qquad e_1^2=e_1,\quad e_2^2=e_2 .
\]

在 $R/I\times R/J$ 中它们正是 $(1,0)$ 与 $(0,1)$ 的原像。

\textbf{于是得到同构。} 核是 $I\cap J$、像满，由第一同构定理

\[
R/(I\cap J)\;\cong\;R/I\times R/J.
\]

有限个两两互素理想同样成立：

\[
R/\bigcap_{k=1}^n I_k\;\cong\;\prod_{k=1}^n R/I_k.
\]

\textbf{一个具体例子。} 取 $R=\mathbb Z$、$I=(3)$、$J=(4)$。取 $i=-3\in I$、$j=4\in J$，则 $i+j=1$，两者互素，于是

\[
\mathbb Z/12\mathbb Z\;\cong\;\mathbb Z/3\mathbb Z\times\mathbb Z/4\mathbb Z .
\]

对应的幂等元是 $e_1=j=4$、$e_2=i=-3$。检验：$4\equiv1\pmod3$、$4\equiv0\pmod4$；$-3\equiv0\pmod3$、$-3\equiv1\pmod4$。给定余数对 $(a,b)$，原像就是

\[
r=4a-3b \pmod{12}.
\]

例如 $(a,b)=(1,1)$ 给出 $r=1$，$(a,b)=(0,1)$ 给出 $r=9$，$(a,b)=(1,0)$ 给出 $r=4$。这也解释了为什么同时给定“模 $3$ 的余数”和“模 $4$ 的余数”能唯一确定“模 $12$ 的余数”。在 $\mathbb Z$ 中 $m\mathbb Z+n\mathbb Z=\mathbb Z$ 当且仅当 $\gcd(m,n)=1$，所以整数版中国剩余定理正是这里的特例。

\textbf{不含幺时 $I\cap J=IJ$ 会失效。} 前面用 $a=a(i+j)=ai+aj$ 推 $I\cap J\subseteq IJ$ 时用到了 $1\in R$。去掉单位元后结论不必成立。取 $R=2\mathbb Z$（不含 $1$ 的环）、$I=J=2\mathbb Z$，则

\[
I\cap J=2\mathbb Z,\qquad IJ=(2\mathbb Z)(2\mathbb Z)=4\mathbb Z,
\]

两者不相等。做这种题时先确认所讨论的环是否含幺元。

\section{单位、分式域与环的特征}


\textbf{本节用到的记号：} $\mathbb Z/n\mathbb Z$ 是模 $n$ 的剩余类环；$R[x]$、$F[x]$ 是多项式环（第 1 节），$F$ 表示域；$\deg f$ 是多项式次数；$\binom pk$ 是二项式系数；$\hookrightarrow$ 表示单同态（嵌入），$\twoheadrightarrow$ 表示满同态；$\exists$ 是“存在”，$\forall$ 是“对所有”。

\subsection{单位与整环}


含幺环 $R$ 的单位群为
\[R^\times=\{u\in R:\exists v\in R,\ uv=vu=1\}.\]
它在乘法下构成群（单位就是“有乘法逆元的元素”）。若存在单位 $u\in R^\times$ 使 $r=us$，就说 $r$ 与 $s$ \textbf{相伴}，记作 $r\approx s$。域满足 $R^\times=R\setminus\{0\}$（每个非零元都可逆）；一般环中还可能有非零但不可逆的元素。

注意本节有两个相似的关系，符号特意分开：\textbf{相伴} $r\approx s$（用一个单位相乘联系起来）与后面分式域里\textbf{分数相等} $(a,b)\sim(c,d)$（交叉相乘相等）。相伴还有另一种写法：在整环中

\[
r\approx s\iff r\mid s\ \text{且}\ s\mid r\iff (r)=(s),
\]

即“互相整除”与“生成同一个主理想”是同一件事。这几种说法在整环中等价，在一般交换环中要小心，本讲只用整环情形。

还有一条常用常识：\textbf{零因子一定不是单位}。因为若 $ab=0$ 且 $a$ 可逆，则 $b=a^{-1}(ab)=0$。所以在 $\mathbb Z/n\mathbb Z$ 中每个元素要么是单位、要么是零因子，这就是 $\mathbb Z/n\mathbb Z$ 是域当且仅当 $n$ 为素数的另一个解释。

整环是交换含幺环、$1\ne0$，且没有非零零因子（与第 1 节的定义一致）。若 $ab=0$ 且 $a\ne0$，则 $b=0$；因此整环中约分成立。$\mathbb Z$ 与域上的多项式环 $F[x]$ 都是整环。

\subsubsection{$F[x]$ 的单位：一个必须写清前提的结论}


课堂上说“$F[x]$ 中的一个多项式可逆当且仅当它的常数项非零”。这句话在 \textbf{$F$ 是域} 时正确，但把它原样推广到任意系数环会出错。正确的分层是：

\begin{itemize}
\item \textbf{$F$ 是域}：$F[x]^\times=F\setminus\{0\}$，即单位恰为非零常数。所以 $1+x$ 不是单位——它的常数项虽然非零，但非常数。
\item \textbf{$R$ 是整环}：$R[x]^\times=R^\times$（[DF] §7.2 Proposition 4）。例如 $\mathbb Z[x]^\times=\{\pm1\}$，而 $2$ 在 $\mathbb Z[x]$ 中不是单位。
\item \textbf{$R$ 是一般交换含幺环}：$f=\sum_{i=0}^n a_ix^i$ 是 $R[x]$ 的单位，当且仅当 \textbf{$a_0$ 是 $R$ 的单位，且 $a_1,\dots,a_n$ 全是幂零元}。这里 $a$ 是\textbf{幂零元}指存在正整数 $m$ 使 $a^m=0$。这正是“常数项要可逆、其余部分要被幂零性吸收”的意思。
\end{itemize}


整环情形可以直接从次数证明：设 $f,g\ne0$ 且 $fg=1$。整环上首项系数之积非零，故

\[
\deg(fg)=\deg f+\deg g.
\]

左边次数为 $0$，所以 $\deg f=\deg g=0$：$f,g$ 都是 $R$ 中的元素，且互为逆元，于是 $f\in R^\times$。

一般环里的幂零条件在 $\mathbb Z/4\mathbb Z$ 上就能看到：$(1+2x)^2=1+4x+4x^2=1$，所以 $1+2x$ 是可逆的，尽管它不是常数。这里 $2$ 是 $\mathbb Z/4\mathbb Z$ 中的幂零元。做多项式环的题时，先确认系数环是哪一类，再决定用哪条结论。

\subsection{分式域}


设 $R$ 是整环。考虑 $R\times(R\setminus\{0\})$ 中的有序对 $(a,b)$，将
\[ (a,b)\sim(c,d)\iff ad=bc \]
视为等价关系，记等价类为 $a/b$。定义
\[\frac ab+\frac cd=\frac{ad+bc}{bd},\qquad \frac ab\frac cd=\frac{ac}{bd}.\]
得到一个域 $\operatorname{Frac}(R)$，其中 $a/b=0$ 当且仅当 $a=0$，且非零元 $a/b$ 的逆为 $b/a$。自然嵌入 $R\hookrightarrow\operatorname{Frac}(R)$ 为 $a\mapsto a/1$。

典型例子是 $\operatorname{Frac}(\mathbb Z)=\mathbb Q$，以及 $\operatorname{Frac}(F[x])=F(x)$，即有理函数域。

\textbf{泛性质（最小性）。} 上面的构造不只是“造出一个域”，它造出的是\textbf{包含 $R$ 的最小域}。精确地说：若 $K$ 是域、$\varphi:R\to K$ 是单同态（即把 $R$ 嵌入 $K$ 的一个同构副本），则存在唯一的域同态 $\widetilde{\varphi}:\operatorname{Frac}(R)\to K$ 使它在 $R$ 上与 $\varphi$ 一致：

\[
\widetilde{\varphi}(a/1)=\varphi(a)\quad(\forall a\in R).
\]

为什么唯一？因为 $\widetilde{\varphi}$ 是域同态，必须满足

\[
\widetilde{\varphi}(a/b)=\widetilde{\varphi}(a/1)\,\widetilde{\varphi}(1/b)=\varphi(a)\,\varphi(b)^{-1},
\]

右边完全由 $\varphi$ 决定，所以至多一种取法。存在性要把这个式子当作定义，再验证它与分数代表元的选取无关——这一步用的正是等价关系 $ad=bc$：若 $a/b=c/d$，即 $ad=bc$，两边施加 $\widetilde{\varphi}$ 即得 $\widetilde{\varphi}(a/b)=\widetilde{\varphi}(c/d)$。

\textbf{这是“最小域”的确切含义：} 任何包含 $R$ 的域 $K$ 都必须含有所有 $\varphi(a)\varphi(b)^{-1}$，也就是必须含有 $\operatorname{Frac}(R)$ 的一个同构副本；所以 $\operatorname{Frac}(R)$ 是包含 $R$ 的域中最小的一个。这也解释了 $\operatorname{Frac}(\mathbb Z)=\mathbb Q$ 的含义：$\mathbb Q$ 是包含 $\mathbb Z$ 的最小域，任何包含 $\mathbb Z$ 的域都含有 $\mathbb Q$。

\subsection{特征}


含幺环的特征定义为
\[\operatorname{char}(R)=\min\{n\ge1:n\cdot1_R=0\},\]
若不存在这样的正整数，则特征为 $0$。若 $R$ 是整环且特征非零，它必为素数：若 $n=ab$ 为合数，则 $(a1_R)(b1_R)=n1_R=0$；无零因子迫使 $a1_R=0$ 或 $b1_R=0$，与最小性矛盾。对一般含幺环，特征可以是合数，例如 $\mathbb Z/4\mathbb Z$ 的特征为 $4$。

特征为素数 $p$ 时有一个非常好用的推论：\textbf{若 $R$ 是特征 $p$ 的交换环，则对任意 $a,b\in R$}

\[
(a+b)^p=a^p+b^p .
\]

这是因为二项式展开 $(a+b)^p=\sum_{k=0}^p\binom pk a^kb^{p-k}$ 中，中间系数 $\binom pk$（$0<k<p$）都被 $p$ 整除，而它们在 $R$ 中作用为乘以 $p\cdot1_R=0$，于是中间项全部消失。有限域理论里的 Frobenius 映射 $x\mapsto x^p$ 就是这个等式。

对每个素数 $p$，$\mathbb F_p=\mathbb Z/p\mathbb Z$ 是含 $p$ 个元素的域，特征为 $p$。若 $R$ 是整环，则存在唯一的单位保持单同态

\[
\mathbb Z/\operatorname{char}(R)\mathbb Z\;\longrightarrow\;R,
\]

它在 $\operatorname{char}(R)=p$ 时给出 $\mathbb F_p\hookrightarrow R$，在 $\operatorname{char}(R)=0$ 时给出 $\mathbb Z\hookrightarrow R$。所以每个整环都含有一个“最小”的子环，称为\textbf{素子环}：特征为 $p$ 时是 $\mathbb F_p$，特征为 $0$ 时是 $\mathbb Z$。注意这里强调整环：一般含幺环上这个映射未必单射（例如 $\mathbb Z/4\mathbb Z$ 上的 $\mathbb Z\to\mathbb Z/4\mathbb Z$ 核为 $4\mathbb Z$）。

因此可以在 $\mathbb F_p[x]$ 上取多项式环，得到特征为 $p$ 的无限整环；再由分式域构造得到特征为 $p$ 的无限域 $\mathbb F_p(x)$。

\subsection{本讲检查清单}


1. \textbf{先确认约定}：本讲的环默认含幺且交换性随语境说明；谈“理想”时默认双边，谈“同态”时默认保幺。
2. \textbf{想证明某物是理想}，要验证非空、对减法封闭、并对\textbf{整个环}的左乘右乘封闭，而不是只对自身元素封闭。
3. \textbf{看到 $R/I$ 先确认 $I$ 是理想}，否则乘法根本不良定义。
4. \textbf{要证 $R/M$ 是域}，用理想对应定理把“$M$ 极大”翻译成“$R/M$ 没有非平凡理想”。
5. \textbf{要证 $R/P$ 是整环}，把“$ab\in P\Rightarrow a\in P$ 或 $b\in P$”翻译成“$R/P$ 无零因子”。
6. \textbf{中国剩余定理}的关键是 $I+J=R$ 给出满射；核总是 $I\cap J$，不需要互素。
7. \textbf{多项式环的单位}先看系数环是哪一类：域 → 非零常数；整环 → $R^\times$；一般环 → 常数项可逆 + 其余系数幂零。
8. \textbf{特征是否必为素数}只对整环成立；一般含幺环可以是合数。

\subsection{参考对照}


\begin{itemize}
\item [DF] §7.1--7.2（环、整环、单位与 $R[x]^\times=R^\times$）、§7.5（环的特征与素子环）、§9.3（分式域的一般构造）。
\item [DS] §10.1（整环与有限整环是域）、§10.3（特征与有限域）。
\item 原始材料：本地归档（手写原稿与录音转写），不纳入课程仓库。
\end{itemize}

\section*{一页式总览}
\begin{itemize}
\item 环：$(R,+)$ 阿贝尔群 $+$ 乘法结合 $+$ 分配律；本讲默认含幺，交换的除环即域。
\item 理想：对加法封闭且对整个环的左右乘法封闭的加法子群；同态的核必是理想，反之每个理想都是某个商映射的核。
\item 第一同构定理：$R/\ker f\cong\operatorname{im}f$；第二：$A/(A\cap I)\cong(A+I)/I$；第三：$(R/I)/(J/I)\cong R/J$。
\item 极大理想 $\iff R/M$ 是域；素理想 $\iff R/P$ 是整环；极大必素，反之不然。
\item 中国剩余定理：$I+J=R$ 时 $R/(I\cap J)\cong R/I\times R/J$，满射性正来自 $I+J=R$。
\item 单位群 $R^\times$；整环上 $R[x]^\times=R^\times$，域上 $F[x]^\times=F\setminus\{0\}$。
\item 分式域 $\operatorname{Frac}(R)$：整环嵌入其中的最小域，$\operatorname{Frac}(\mathbb Z)=\mathbb Q$。
\item 特征：整环的非零特征必为素数；一般含幺环可以是合数。
\end{itemize}

\section*{参考资料}
\begin{itemize}
\item \textbf{[DF]} D. S. Dummit, R. M. Foote. \emph{Abstract Algebra}, 3rd ed. §§7.1--7.5（环、多项式环、商环、特征与素子环）、§9.3（分式域的一般构造）。
\item \textbf{[DS]} 邓小铁、王畅。《离散数学与结构》§§10.1--10.3（环和域、同态与理想、有限域）；第二、第三同构定理在 §10.2 留作习题。
\item 原始材料：\texttt{DiscreteCourse2026/课堂手记和录音转写/Lec09/}。
\end{itemize}
\end{document}
